% Figure: Lorenz curve and the Gini coefficient as an integral. The
% Lorenz curve L(p) plots the cumulative share of a quantity held by the
% bottom fraction p of the population; the area between the line of
% equality and the Lorenz curve, normalised, is the Gini coefficient.
\begin{figure}[htbp]
\centering
\begin{tikzpicture}
\begin{axis}[
  width=8.8cm, height=8.2cm,
  xlabel={population fraction $p$}, ylabel={cumulative share $L(p)$},
  xmin=0, xmax=1.0, ymin=0, ymax=1.0,
  domain=0:1, samples=120,
  axis lines=left, font=\small,
  axis equal image,
  legend pos=north west, legend style={font=\footnotesize},
]
% line of equality
\addplot[engblack, thick] coordinates {(0,0) (1,1)};
\addlegendentry{line of equality $L=p$}
% Lorenz curve L(p) = p^3 (a convex power-law example)
\addplot[engblue, very thick] {x^3};
\addlegendentry{Lorenz curve $L(p)$}
% shade the gap between equality line and Lorenz curve
\addplot[draw=none, fill=engamber!22, domain=0:1]
  {x} \closedcycle;
\addplot[draw=none, fill=white, domain=0:1]
  {x^3} \closedcycle;
% redraw curves over the white fill
\addplot[engblack, thick] coordinates {(0,0) (1,1)};
\addplot[engblue, very thick] {x^3};
\node[font=\scriptsize, engamber, anchor=center]
  at (axis cs:0.62,0.42) {area $=G/2$};
\end{axis}
\end{tikzpicture}
\caption[Lorenz curve and the Gini coefficient]{The Lorenz curve $L(p)$
gives the cumulative share of a quantity (income, energy use, packet
traffic) held by the poorest fraction $p$ of a population. Perfect
equality is the diagonal $L(p)=p$; any real distribution bows below it.
The amber area between the diagonal and the curve, taken as twice itself
so the maximum equals one, is the Gini coefficient $G = 1 - 2\int_0^1
L(p)\,\dd p$. The curve shown, $L(p)=p^3$, gives $\int_0^1 p^3\,\dd p =
\tfrac14$ and $G = 1 - \tfrac12 = \tfrac12$. The Gini coefficient is a
single integral that compresses a whole distribution into one inequality
number, and the same construction measures load imbalance across servers
or wear across a bearing race.}
\label{fig:vol02ch06:lorenz-gini}
\end{figure}
